\bookchapter{Capítulo 2: Ordena tu habitación digital — Gestión de archivos}

\chapterintro{Con el uso prolongado, el teléfono se desordena; los buenos hábitos empiezan por la gestión de archivos. Igual que tu habitación necesita ordenarse periódicamente, tu espacio digital también requiere mantenimiento. Un teléfono bien organizado no solo facilita encontrar cosas, sino que también funciona con mayor fluidez.}

\booksection{I. Sistema de almacenamiento de tres niveles: tu almacén digital}

\begin{enumerate}
  \item \textbf{Almacenamiento interno}: zona central, alta velocidad, aquí están el sistema y las App. Es el espacio más valioso de tu teléfono, hay que administrarlo con cuidado.
  \item \textbf{Tarjeta SD (si es compatible)}: armario ampliable, menor velocidad, ideal para archivos grandes como películas y fotos. Nota: muchos teléfonos de gama alta ya no admiten tarjetas SD.
  \item \textbf{Almacenamiento en la nube}: caja fuerte + disco en red sincronizado. Previene pérdidas y permite sincronizar entre varios dispositivos. iCloud, Baidu Netdisk y OneDrive son opciones habituales.
\end{enumerate}

\booksection{II. Conoce el “Administrador de archivos”}

Android suele tener una App [Archivos] o [Administrador de archivos]; iOS incluye la App [Archivos]. Son el “explorador” con el que puedes ver todos los archivos de tu teléfono.

\booksubsection{Significado de las carpetas comunes:}

\begin{itemize}
  \item \textbf{DCIM}: fotos y videos tomados con la cámara.
  \item \textbf{Pictures}: otras imágenes, como capturas de pantalla e imágenes guardadas.
  \item \textbf{Download}: archivos descargados por el navegador.
  \item \textbf{Music}: archivos de música.
  \item \textbf{Movies}: archivos de video.
  \item \textbf{Documents}: archivos de documentos.
  \item \textbf{Android}: carpeta del sistema Android, contiene los datos de cada App; por lo general no debe eliminarse manualmente.
\end{itemize}

\booksection{III. Estructura del sistema de archivos: profundización}

Android se basa en Linux y su sistema de archivos tiene cierta estructura jerárquica. Aunque el usuario común no necesita conocerla a fondo, entender algúnos conceptos básicos resulta útil.

\booksubsection{Particiones de almacenamiento:}

\begin{itemize}
  \item \textbf{/system}: partición del sistema, almacena los archivos del sistema operativo. El usuario común no tiene acceso; requiere permisos root.
  \item \textbf{/data}: partición de datos, almacena los datos de usuario y de aplicaciones. Es el espacio que usas a diario.
  \item \textbf{/sdcard o /storage/emulated/0}: espacio de almacenamiento accesible para el usuario; aquí están tus fotos y archivos descargados.
\end{itemize}

\booksubsection{Almacenamiento de datos de App:}

\begin{itemize}
  \item \textbf{Almacenamiento interno}: espacio privado de la App, otras App no pueden acceder. Se elimina al desinstalar la App.
  \item \textbf{Almacenamiento externo}: archivos que la App guarda en la zona accesible para el usuario, como fotos y descargas. No siempre se eliminan al desinstalar la App.
\end{itemize}

\booksection{IV. Copia de seguridad y restauración de datos: asegura tus datos}

Los datos no tienen precio; hacer copias de seguridad es un hábito imprescindible.

\booksubsection{Estrategias de copia de seguridad:}

\begin{enumerate}
  \item \textbf{Copia en la nube}: automática y cómoda, pero requiere red y espacio en la nube. iCloud, Huawei Cloud, Xiaomi Cloud, entre otros, son compatibles.
  \item \textbf{Copia local}: conecta el teléfono a la computadora con un cable y usa la herramienta oficial. Rápida y de gran capacidad, pero requiere operación manual.
  \item \textbf{Copia doble}: copia en la nube + copia local, la más segura.
\end{enumerate}

\booksubsection{¿Qué respaldar?}

\begin{itemize}
  \item \textbf{Fotos y videos}: lo que más espacio ocupa y lo más valioso.
  \item \textbf{Contactos}: lo más cómodo es guardarlos en la cuenta de la nube.
  \item \textbf{Historial de chat de WeChat}: usa la función de copia de seguridad integrada de WeChat.
  \item \textbf{Documentos importantes}: guárdalos en la nube o en la computadora.
\end{itemize}

\booksection{V. Seguridad de archivos: protege tu privacidad}

\begin{itemize}
  \item \textbf{Cifrado de archivos sensibles}: muchos teléfonos tienen una función de "carpeta segura" o "espacio privado" donde puedes guardar archivos sensibles; requieren contraseña o huella para acceder.
  \item \textbf{Eliminar no equivale a destruir}: la eliminación común solo marca el espacio como reutilizable, los datos siguen ahí. Para borrarlos por completo, usa herramientas de borrado seguro especializadas o sobrescribe varias veces.
  \item \textbf{Seguridad en la nube}: elige proveedores de nube de buena reputación, activa la verificación en dos pasos y revisa periódicamente los dispositivos conectados.
\end{itemize}

\booksection{VI. ¿Sin espacio? Haz una limpieza a fondo}

\begin{itemize}
  \item \textbf{Limpiar la caché}: la memoria temporal de las App; al limpiarla liberas espacio. Se puede hacer en Configuración > Almacenamiento.
  \item \textbf{Buscar archivos grandes}: usa el administrador de archivos para ordenar por tamaño y encontrar a los "culpables" que ocupan espacio.
  \item \textbf{Eliminar restos}: tras desinstalar una App, a veces quedan archivos basura; elimínalos manualmente.
  \item \textbf{Limpiar el álbum}: borra fotos duplicadas, fotos borrosas y capturas de pantalla innecesarias.
  \item \textbf{Limpiar WeChat}: WeChat es un gran consumidor de espacio; es necesario limpiar periódicamente el historial de chat y la caché.
\end{itemize}

\booksection{VII. Análisis completo de los formatos de archivo comunes}

Las extensiones de archivo en el teléfono pueden abrumar; entenderlas te ayudará a gestionar mejor los archivos:

\begin{center}
\begin{tabularx}{\linewidth}{|X|X|X|X|}
\hline
\textbf{Formato} & \textbf{Tipo} & \textbf{Características} & \textbf{Apertura recomendada} \tabularnewline
\hline
JPG/JPEG & Imagen & Compresión con pérdida, tamaño pequeño, máxima compatibilidad. Formato predeterminado de fotos del teléfono. & Álbum/cualquier visor de imágenes \tabularnewline
PNG & Imagen & Compresión sin pérdida, admite fondo transparente. Común en capturas de pantalla e iconos. Mayor tamaño que JPG. & Álbum/cualquier visor de imágenes \tabularnewline
MP3 & Audio & Compresión con pérdida, tamaño pequeño, excelente compatibilidad. El formato de audio más universal. & Cualquier reproductor de música \tabularnewline
FLAC & Audio & Compresión sin pérdida, calidad idéntica al CD, tamaño 3-5 veces mayor que MP3. El favorito de los audiófilos. & Reproductor de música/App reproductor profesional \tabularnewline
WAV & Audio & Audio original sin comprimir, la mejor calidad pero tamaño enorme. Común en grabaciones profesionales. & Reproductor de música \tabularnewline
M4A & Audio & Codificación AAC, mejor calidad que MP3 a la misma tasa de bits. Común en el ecosistema Apple. & Apple Music/reproductor de música \tabularnewline
MP4 & Video & El formato de contenedor de video más universal, excelente compatibilidad. Formato predeterminado de los videos grabados por el teléfono. & Álbum/cualquier reproductor de video \tabularnewline
MKV & Video & Formato de contenedor de código abierto, admite múltiples pistas de audio y subtítulos. Común en recursos de películas. & VLC/nPlayer \tabularnewline
DOC/DOCX & Documento & Documento Word. DOCX es un formato comprimido, de menor tamaño. DOC es el formato antiguo. & WPS/Office \tabularnewline
PPT/PPTX & Presentación & Presentación de PowerPoint. PPTX es un formato comprimido. & WPS/Office \tabularnewline
XLS/XLSX & Hoja de cálculo & Hoja de cálculo Excel. XLSX es un formato comprimido. & WPS/Office \tabularnewline
PDF & Documento & Formato multiplataforma, diseño fijo no editable. Común en contratos, tesis y libros electrónicos. & Adobe Reader/WPS \tabularnewline
ZIP & Compresión & El formato de compresión más universal, buena compatibilidad multiplataforma. & Sistema integrado/App de descompresión \tabularnewline
RAR & Compresión & Mayor tasa de compresión que ZIP, pero requiere software especializado para descomprimir. & App de descompresión \tabularnewline
ISO & Imagen & Archivo de imagen de disco, copia completa del contenido de un disco. Común en paquetes de instalación de sistemas. & App de descompresión/unidad virtual \tabularnewline
HTML & Página web & Archivo de página web, se abre con el navegador. Las páginas web guardadas suelen usar este formato. & Navegador \tabularnewline
\hline
\end{tabularx}
\end{center}

\begin{tcbpractice}{Practica: ordena la carpeta "Descargas"}
Abre el administrador de archivos > entra en la carpeta “Descargas”. Mira qué hay dentro. Elimina lo inútil y clasifica lo útil (por ejemplo, crea carpetas como “Materiales de trabajo” o “Stickers” y guárdalos allí). Aprovecha para revisar el espacio de almacenamiento y saber qué tipos de archivo ocupan más.

\end{tcbpractice}

\booksection{VIII. Gestión de imágenes y sistema de archivos: mantén tus activos en orden}

En la sección anterior aprendimos a “conocer las carpetas”, pero conocer no es suficiente. Esta sección trata sobre\textbf{cómo establecer activamente un conjunto de normas de gestión aplicables a largo plazo}, para que tus archivos siempre se puedan encontrar, no se pierdan ni sean espiados.

\booksubsection{Normas de clasificación y nomenclatura de archivos}

La raíz del caos es la ausencia de reglas. Establecerlas es sencillo, la clave está en\textbf{la constancia}.

\begin{itemize}
  \item \textbf{Nomenclatura "fecha + evento"}: incluye en el nombre de la carpeta una marca de tiempo y una descripción del evento, para facilitar el recuerdo y la búsqueda. Por ejemplo,\texttt{202606\_RecibosViaje},\texttt{2025ResumenAnual},\texttt{2024CumpleañosNiño}. Así puedes ver la línea temporal y conocer el contenido de un vistazo.
  \item \textbf{Sistema de clasificación de tres niveles}:
  \begin{itemize}
    \item \textbf{Nivel 1}: por tipo de archivo (documentos / imágenes / videos / audio)
    \item \textbf{Nivel 2}: por escenario de uso (trabajo / vida personal / estudio)
    \item \textbf{Nivel 3}: por evento o proyecto concreto (por ejemplo, "2026Q2\_ProyectoA")
  \end{itemize}
  \item \textbf{Plan mensual de limpieza}: elige una fecha fija cada mes (se recomienda el día de cobro o el día 1 del mes) para organizar el directorío "Descargas":\textbf{archiva lo útil, elimina lo inútil, guarda lo valioso}. Doce pequeñas limpiezas al año son mucho más ligeras que una gran limpieza anual.
\end{itemize}

\begin{tcbpractice}{Lista de tareas de este mes}
Elige un fin de semana libre y sigue este orden:

\begin{enumerate}
  \item Crea en el administrador de archivos una carpeta de trabajo de este mes con el formato\texttt{2026Q2\_ProyectoA};
  \item Revisa todos los archivos de esta semana en el directorío "Descargas", conserva lo útil, archívalo y elimina lo inútil;
  \item Crea carpetas anuales para las fotos familiares de este año, como\texttt{2025Familia},\texttt{2026Familia}.
\end{enumerate}

\end{tcbpractice}

\booksubsection{Estrategia para archivos redundantes}

Con el uso prolongado, el teléfono “engorda”; la culpa suele ser de los archivos redundantes.

\begin{itemize}
  \item \textbf{Limpieza automática a las 72 horas}: los archivos temporales recibidos en App sociales como WeChat y QQ (stickers, imágenes de chat, videos cortos) suelen ser inútiles tras una primera vista. En [WeChat] > [Configuración] > [General] > [Espacio de almacenamiento], activa la "limpieza automática" y ajusta la caché caducada a 72 horas. El teléfono quedará mucho más ligero al instante, sin perder registros importantes.
  \item \textbf{Limpieza dirigida del espacio de almacenamiento}: en Android [Configuración] > [Almacenamiento] > [Limpieza y optimización]; en iOS [Configuración] > [General] > [Espacio de almacenamiento del iPhone], el sistema clasificará los elementos en "caché / archivos temporales / App desinstalables / archivos grandes". Concéntrate en: (1) archivos de caché; (2) paquetes de instalación caducados (los paquetes de versiones antiguas que quedan tras actualizar las App); (3) archivos duplicados (fotos con el mismo nombre y tamaño).
  \item \textbf{Reglas de clasificación automática}: la mayoría de las App de álbum de las marcas admiten la clasificación inteligente: personas, paisajes, documentos, capturas de pantalla. Al activarla, una foto grupal nueva se archivará automáticamente en "Personas > Zhang San", y un contrato escaneado en "Documentos".\textbf{Combinado con la organización mensual, reduce drásticamente la probabilidad de "no encontrar las cosas"}.
\end{itemize}

\booksubsection{Normas de tratamiento seguro de imágenes}

Las imágenes son el tipo de dato más “sensible” del teléfono: una foto casual puede revelar tu ubicación, dispositivo y hora. Recomendamos\textbf{establecer un flujo de trabajo de “doble protección”}:

\begin{tcbwarn}{Procedimiento estándar para enviar fotos de identificación o capturas sensibles}
\begin{enumerate}
  \item \textbf{Difuminar la información visible}: con herramientas del sistema o de terceros (como Meitu XiuXiu, Snapseed)\textbf{cifra los dígitos centrales del número de identificación, las matrículas, los números de puerta, los nombres y los importes};
  \item \textbf{Eliminar metadatos}: antes de enviar, en el álbum pulsa "Compartir" o "Exportar"; el sistema preguntará si deseas conservar la información de ubicación, elige\textbf{"No conservar"}. Así, el modelo del dispositivo, las coordenadas GPS y la hora de toma no se enviarán junto con la imagen;
  \item \textbf{Marca de agua temporal}(opcional): añade a las capturas sensibles una marca de agua de texto tipo "Solo para trámite XX", para evitar usos indebidos.
\end{enumerate}

\end{tcbwarn}

Métodos comunes para limpiar metadatos EXIF (operaciones básicas; para protección profunda, ver Capítulo 6 “Anonimización de fotos”):

\begin{itemize}
  \item \textbf{Sistema integrado}: en iOS [Configuración] > [Privacidad y seguridad] > [Cámara], puedes desactivar "Ubicación precisa"; al compartir, selecciona "Solo esta foto" y\textbf{se eliminarán automáticamente la mayoría de los EXIF}.
  \item \textbf{Android}: la mayoría de las App de álbum de las marcas preguntan al compartir si deseas "eliminar la información de ubicación"; márcalo.
  \item \textbf{Herramientas profesionales}: ExifTool (en computadora), Metapho (iOS), Photo Metadata Remover (Android) pueden eliminar por completo todos los metadatos.
\end{itemize}

Para imágenes sensibles que necesites\textbf{conservar a largo plazo}(identificación, tarjetas bancarias, escaneos del registro civil, pasaporte), se recomienda guardarlas en un\textbf{álbum cifrado}:

\begin{itemize}
  \item iPhone: [Fotos] > selecciona la foto > Ocultar (solo visible tras verificar Face ID);
  \item Huawei/Xiaomi: [Álbum] > mantén pulsada la foto > "Mover a la caja fuerte";
  \item App de cifrado de terceros:\textbf{360 Álbum seguro},\textbf{Administrador de álbumes de Tencent (álbum privado)},\textbf{KeepSafe}, entre otros, con protección por contraseña independiente y mayor discreción.
\end{itemize}
